% Arabic Sentiment Analysis — Report structure (from docs/repport-struct.md)
% Compile: pdflatex report-from-struct.tex (twice for TOC). For Arabic examples in the PDF body, use XeLaTeX + fontspec + an Arabic font.

\documentclass[11pt,a4paper]{report}

\usepackage[a4paper,margin=2.5cm]{geometry}
\usepackage[T1]{fontenc}
\usepackage[utf8]{inputenc}
\usepackage{amsmath,amssymb}
\usepackage{graphicx}
\usepackage{booktabs}
\usepackage[hidelinks]{hyperref}
\usepackage{enumitem}
\usepackage{setspace}
\usepackage{xcolor}
\usepackage{listings}
\lstdefinestyle{py}{
  language=Python,
  basicstyle=\ttfamily\footnotesize,
  keywordstyle=\color{blue},
  commentstyle=\color{gray},
  stringstyle=\color{teal},
  breaklines=true,
  showstringspaces=false,
  columns=fullflexible,
}
\usepackage{fancyvrb}

\title{Sentiment Analysis in Arabic:\\Comparative Study of ML and DL Approaches}
\author{Your Name(s)\\\small Module name --- Faculty / Department\\\small Supervised by: Professor Name}
\date{Academic Year 2025--2026}

\begin{document}

\maketitle

\begin{center}
\textbf{University Name}\\[0.5em]
Faculty / Department\\
\end{center}

\thispagestyle{empty}
\vfill
\begin{center}
\textit{Replace placeholders above with your institution details.}
\end{center}
\newpage

\tableofcontents
\clearpage

% --- 1 ---
\chapter{Introduction}

\section{Goal}
Explain what sentiment analysis is, why Arabic sentiment analysis is challenging, and what this project studies.

\subsection*{Motivation (example points)}
\begin{itemize}[noitemsep]
    \item Growth of Arabic social media content.
    \item Importance of NLP for Arabic.
    \item Challenges specific to Arabic: dialects, morphology, informal writing.
\end{itemize}

\section{Objectives}
State your objectives explicitly. Example:

\textit{This project aims to compare several machine learning and deep learning approaches for Arabic sentiment classification using different text representations and feature selection techniques.}

% --- 2 ---
\chapter{Related Work}
Brief survey (recommended; keep to 1--2 pages):
\begin{itemize}[noitemsep]
    \item Prior Arabic sentiment analysis research.
    \item Common techniques (e.g., TF-IDF + SVM).
    \item Role of pretrained models such as AraBERT.
\end{itemize}
Add citations (AraBERT papers, TF-IDF + SVM studies, etc.).

% --- 3 ---
\chapter{Dataset Description}

This project uses two complementary Arabic corpora stored under \texttt{data/}: a microblog-style tweet collection (\texttt{Tweets.txt}) and a large tabular file of hotel guest reviews (\texttt{balanced-reviews.txt}). Together they cover \emph{short, noisy social text} versus \emph{longer, structured review text}, which is useful when comparing machine learning and deep learning models and when discussing domain shift.

\begin{spacing}{1.1}

\section{Arabic tweets corpus (\texttt{data/Tweets.txt})}

\subsection{Origin and format}
The file \texttt{Tweets.txt} contains one record per line in the form \texttt{text\textbackslash tLABEL}. The tweet appears first; the sentiment or stance label is separated by a tab and appears at the end of the line. All parsing and statistics below assume UTF-8 encoding. Blank lines are omitted; after cleaning empty rows, the corpus contains \textbf{10\,006} labeled tweets.

\subsection{Labels and interpretation}
The scheme uses \textbf{four} labels (uppercase in file):
\begin{itemize}[noitemsep]
    \item \textbf{POS} --- subjective positive sentiment (praise, congratulations, admiration).
    \item \textbf{NEG} --- subjective negative sentiment (criticism, rejection, complaint).
    \item \textbf{NEUTRAL} --- factual or mixed tone without a clear positive/negative stance toward a single target.
    \item \textbf{OBJ} --- objective or informational content (news-style headlines, event announcements, invitations), often without evaluative opinion.
\end{itemize}
This is characteristic of Arabic Twitter datasets where ``objective'' and ``neutral'' are kept distinct: many posts are purely informational (\textbf{OBJ}), which differs from binary positive/negative movie or product reviews.

\subsection{Class distribution and implications}
The empirical distribution over the 10\,006 tweets is:
\begin{center}
\begin{tabular}{@{}lrr@{}}
    \toprule
    \textbf{Label} & \textbf{Count} & \textbf{Share (\%)} \\
    \midrule
    OBJ     & 6\,691 & $\approx$66.9 \\
    NEG     & 1\,684 & $\approx$16.8 \\
    NEUTRAL &   832 & $\approx$8.3 \\
    POS     &   799 & $\approx$8.0 \\
    \midrule
    \textbf{Total} & 10\,006 & 100 \\
    \bottomrule
\end{tabular}
\end{center}
The corpus is \textbf{strongly imbalanced}: objective tweets dominate. Any model trained without adjustment may achieve high accuracy by predicting OBJ most of the time. For fair comparison, macro-averaged F1, per-class recall, or class-weighted training should be reported. The text itself is typical of Twitter: hashtags, mentions, code-switching, dialectal Egyptian and Levantine forms, elongation, and occasional non-Arabic tokens appear---so preprocessing (normalization, handling hashtags) materially affects results.

\subsection{Typical content (qualitative)}
Posts range from short reactions and political commentary to event invitations and news snippets. In the current file, tweet bodies average about \textbf{83} characters (maximum about 178 before the tab label). This dataset stresses \emph{noise robustness} and \emph{multi-class} learning rather than long-form argumentation.

\section{Balanced hotel reviews (\texttt{data/balanced-reviews.txt})}

\subsection{Origin and format}
The file \texttt{balanced-reviews.txt} is \textbf{UTF-16 LE} (Unicode BOM \texttt{FF FE}); it must be read with the correct encoding when loading in Python or other tools. Each logical row is tab-separated with seven fields:
\begin{enumerate}[noitemsep]
    \item Row index (\texttt{no}),
    \item Hotel name (Arabic),
    \item Numerical \textbf{rating} on a 1--5 star scale,
    \item Guest type (e.g.\ solo, couple, family),
    \item Room type or ``\texttt{-}'',
    \item Stay duration phrase (e.g.\ ``one night''),
    \item Free-text \textbf{review} in Arabic (possibly with quotes and punctuation).
\end{enumerate}
Using a robust CSV/TSV reader, \textbf{105\,698} complete review rows are obtained (header excluded). The reviews refer to \textbf{1\,173} distinct hotel names, i.e.\ a multi-property hospitality domain comparable in spirit to review datasets used in Arabic sentiment benchmarks (e.g.\ HARD-style accommodation feedback).

\subsection{Star ratings}
The star column is intended as a coarse sentiment proxy. In this file, only ratings \textbf{1}, \textbf{2}, \textbf{4}, and \textbf{5} occur; there is \textbf{no} rating~3 in the extracted rows (likely removed during balancing or filtering). Counts are:
\begin{center}
\begin{tabular}{@{}lrr@{}}
    \toprule
    \textbf{Rating} & \textbf{Count} & \textbf{Share (\%)} \\
    \midrule
    1 (worst)  & 14\,382 & $\approx$13.6 \\
    2          & 38\,467 & $\approx$36.4 \\
    4          & 26\,450 & $\approx$25.0 \\
    5 (best)   & 26\,399 & $\approx$25.0 \\
    \midrule
    \textbf{Total} & 105\,698 & 100 \\
    \bottomrule
\end{tabular}
\end{center}
For binary sentiment experiments, a common mapping is: ratings 4--5 $\rightarrow$ positive, 1--2 $\rightarrow$ negative (rating~3 would be neutral if present). The intermediate rating~2 is frequent, so the corpus contains many \emph{weakly negative} or mixed judgments expressed in full sentences.

\subsection{Linguistic characteristics}
Reviews are substantially \textbf{longer} than tweets: the review field averages about \textbf{137} characters per row in this file, with a long tail (maximum about 3\,365 characters for a single review). Language is often Modern Standard Arabic mixed with Gulf or Levantine usage; guests describe cleanliness, staff, location, breakfast, and value. Compared to \texttt{Tweets.txt}, vocabulary is more repetitive (hotel-specific collocations) but spelling is usually more standard and less hashtag-driven.

\section{Summary comparison}

\begin{table}[htbp]
    \centering
    \caption{Overview of the two corpora (counts from the project files as analysed)}
    \begin{tabular}{@{}p{3.2cm}p{4.5cm}p{5.8cm}@{}}
        \toprule
        \textbf{Aspect} & \textbf{\texttt{Tweets.txt}} & \textbf{\texttt{balanced-reviews.txt}} \\
        \midrule
        Size & 10\,006 tweets & 105\,698 reviews \\
        Labels & 4-class: POS, NEG, NEUTRAL, OBJ & 1--5 stars (1,2,4,5 present) \\
        Encoding & UTF-8 & UTF-16 LE \\
        Domain & Social media, political/cultural mix & Hotel stays (multi-hotel) \\
        Length & Short & Longer paragraphs \\
        Balance & Highly skewed toward OBJ & Rating frequencies deliberately structured \\
        \bottomrule
    \end{tabular}
\end{table}

\noindent
These differences motivate separate preprocessing choices (e.g.\ hashtag handling for tweets, sentence segmentation for reviews) and careful interpretation of metrics when aggregating or comparing results across corpora.

\end{spacing}

% --- 4 ---
\chapter{Data Preprocessing}

Two corpora are cleaned with \textbf{different pipelines} implemented in \texttt{src/preprocessing.py}: \texttt{preprocess\_tweet} for microblog text (\texttt{data/Tweets.txt}) and \texttt{preprocess\_balanced\_review} for hotel feedback (\texttt{data/balanced-reviews.txt}). Quantitative comparisons below were produced by \texttt{scripts/analyze\_preprocessing.py} on the full files (statistics saved as \texttt{report/preprocessing\_stats.json}).

\section{Pipeline for tweets (\texttt{preprocess\_tweet})}

\begin{enumerate}[noitemsep]
    \item \textbf{Social tokens:} remove URLs; drop \texttt{@}mentions; replace hashtags \texttt{\#word} by \texttt{word} so topical Arabic is retained; collapse long underscore/dash runs typical of noisy posts.
    \item \textbf{Script filter:} remove Arabic diacritics not needed for polarity; keep Arabic letters and spaces; squeeze whitespace.
    \item \textbf{Surface normalisation:} optional emoji$\rightarrow$sentiment placeholders; strip Latin punctuation; unify alif/yaa/\textit{taa marbuta} variants (e.g.\ \textit{alef} and \textit{haa} forms).
    \item \textbf{Deep clean:} repeated-character shortening (e.g.\ elongation), tatweel removal, non-Arabic symbol stripping while keeping underscore inside tokens when useful.
\end{enumerate}

Stemming is \textbf{optional} (\texttt{use\_stemming=True}); CSV export in \texttt{src/load\_to\_csv.py} enables it for archived tweet files.

\subsection{Usage (code)}
Listing~\ref{lst:prep-import} shows how the cleaners and binary star mapping are imported and called from application code or notebooks.

\begin{lstlisting}[style=py, caption={Import and single-document cleaning.}, label={lst:prep-import}]
from src.preprocessing import (
    preprocess_tweet,
    preprocess_balanced_review,
    preprocess_tweets,
    preprocess_balanced_reviews,
    star_rating_to_binary,
)

text_t = preprocess_tweet(raw_tweet_string)
text_r = preprocess_balanced_review(raw_review_string)

y_bin = star_rating_to_binary(row_rating)  # "POS", "NEG", or None
\end{lstlisting}

\begin{lstlisting}[style=py, caption={Batch cleaning before vectorisation.}, label={lst:prep-batch}]
X_tweets_clean = preprocess_tweets(list_of_raw_tweets)
X_reviews_clean = preprocess_balanced_reviews(list_of_raw_review_cells)
\end{lstlisting}

\section{Pipeline for hotel reviews (\texttt{preprocess\_balanced\_review})}

\begin{enumerate}[noitemsep]
    \item \textbf{Decorative quotes:} remove smart/ASCII quotes often pasted around review titles in scraped tables.
    \item Same \textbf{script filter}, \textbf{emoji/punctuation/normalisation}, and \textbf{deep clean} as generic Arabic text (hashtag logic is disabled---reviews rarely use Twitter-style \texttt{\#}).
\end{enumerate}

\section{Binary labels for star ratings}

For experiments that need \textbf{two classes} on \texttt{balanced-reviews.txt}, numeric ratings are mapped with \texttt{star\_rating\_to\_binary}: ratings 4--5 $\rightarrow$ \texttt{POS}, 1--2 $\rightarrow$ \texttt{NEG}; rating~3 can be dropped or folded (parameter \texttt{neutral}). The current extract contains only stars $\{1,2,4,5\}$.

\begin{lstlisting}[style=py, caption={Binary mapping for the rating column (strings or ints).}, label={lst:prep-binary}]
label = star_rating_to_binary("5")   # -> "POS"
label = star_rating_to_binary(1)    # -> "NEG"
label = star_rating_to_binary(3)      # -> None if neutral="drop" (default)
\end{lstlisting}

\section{Illustrative text transformations}

Real lines were taken from \texttt{data/Tweets.txt} and \texttt{data/balanced-reviews.txt} and processed with the project pipelines (see \texttt{scripts/sample\_preprocessing\_examples.py}). Table~\ref{tab:text-transform} summarises the \emph{linguistic} effect; the exact UTF-8 strings are reproduced below.

\begin{table}[htbp]
    \centering
    \caption{Qualitative summary of two sample transformations.}
    \label{tab:text-transform}
    \begin{tabular}{@{}p{2.2cm}p{11cm}@{}}
        \toprule
        \textbf{Corpus} & \textbf{What changes} \\
        \midrule
        Tweet &
        Hashtag markers (\texttt{\#}) are removed while keeping the following Arabic letters as ordinary tokens; URLs and \texttt{@}mentions are removed; common orthographic rules merge alef variants and map \textit{ta marbuta} to \textit{haa} where configured. \\
        \midrule
        Review &
        Decorative quotes and stray punctuation copied from booking tables are stripped before the same Arabic normalisation pass used for tweets (reviews do not use Twitter hashtag rules). \\
        \bottomrule
    \end{tabular}
\end{table}

\noindent\textit{Note:} Run \LaTeX{} from the \texttt{report/} folder so \verb|preprocessing_examples.utf8.txt| is found. If Arabic glyphs show as blanks, use \textbf{XeLaTeX}/\textbf{LuaLaTeX} with an Arabic font, or read \texttt{preprocessing\_examples.utf8.txt} directly.

\vspace{0.5em}
\small
\VerbatimInput{preprocessing_examples.utf8.txt}
\normalsize

\section{Empirical effect of cleaning (full-corpus statistics)}

Cleaning shortens texts slightly by removing URLs, punctuation, and decorative symbols while preserving lexical content. Table~\ref{tab:preprocess-stats} summarises length and token counts \emph{before} (raw) and \emph{after} (cleaned) preprocessing.

\begin{table}[htbp]
    \centering
    \caption{Raw vs.\ cleaned statistics (character length and whitespace-separated token counts).}
    \label{tab:preprocess-stats}
    \begin{tabular}{@{}lrrrr@{}}
        \toprule
        & \multicolumn{2}{c}{\textbf{Characters}} & \multicolumn{2}{c}{\textbf{Tokens}} \\
        \cmidrule(lr){2-3} \cmidrule(lr){4-5}
        \textbf{Corpus ($n$)} & \textbf{Mean} & \textbf{Median} & \textbf{Mean} & \textbf{Median} \\
        \midrule
        Tweets raw (10\,006) & 83.6 & 86 & 14.2 & 14 \\
        Tweets cleaned & 78.9 & 80 & 14.2 & 14 \\
        \midrule
        Reviews raw (105\,698) & 137.1 & 96 & 24.0 & 16 \\
        Reviews cleaned & 130.0 & 91 & 23.2 & 16 \\
        \bottomrule
    \end{tabular}
\end{table}

\noindent
\textbf{Relative shortening:} mean character length drops by about \textbf{6.4\,\%} for tweets and \textbf{7.0\,\%} for reviews (URLs, punctuation, and quote stripping). Token counts change modestly ($\approx$0.3 tokens fewer per tweet and $\approx$0.8 per review on average), showing that the pipelines mainly \textit{normalise surface form} rather than aggressively truncate content.

\section{Discussion: differences between corpora after cleaning}

\begin{itemize}[noitemsep]
    \item \textbf{Length:} cleaned reviews remain longer on average than tweets ($\approx$130 vs.\ $\approx$79 characters mean), consistent with paragraph-style hotel feedback versus short posts.
    \item \textbf{Variance:} reviews show higher standard deviation in length (long outliers exist); tweets are more tightly clustered around a short length band.
    \item \textbf{Empty strings:} a handful of rows become empty after cleaning (5 tweets, 67 reviews in the measured run), mostly noise-only lines; these rows should be dropped before training.
    \item \textbf{Design implication:} hashtag-preserving cleaning suits opinion and topic cues in Twitter data; quote stripping suits tabular review exports. Using one generic cleaner for both domains would either lose hashtag semantics or skip hotel-specific noise removal.
\end{itemize}

\section{Optional stemming}

If enabled, a light Arabic stemmer (package \texttt{light\_stemmer} when installed, else a rule-based fallback) reduces morphological variants to a common surface form; enable only when evaluating sparse bag-of-words models, since stemming can hurt neural models that rely on full word shapes.

% --- 5 ---
\chapter{Text Representation}

Machine learning and many deep models cannot consume raw strings: each document must be mapped to a fixed-length numerical representation (sparse counts, dense vectors, or contextual embeddings). This chapter summarises the representations used in Arabic sentiment analysis and relates them to the implementations in \texttt{src/Feature\_Extraction.py} and \texttt{src/embeddings.py}.

\section{Problem setting}

Let $\mathcal{D} = \{d_1,\ldots,d_N\}$ be a corpus of preprocessed Arabic texts. A vocabulary $\mathcal{V}$ of terms (words or \textit{n}-grams) is built from training data. Each document $d$ is mapped to a vector $\mathbf{x}(d) \in \mathbb{R}^{|\mathcal{V}|}$ (sparse) or $\mathbb{R}^{k}$ with $k \ll |\mathcal{V}|$ (dense), or to a sequence of token IDs for neural encoders.

\begin{table}[htbp]
    \centering
    \caption{Overview of representation families (typical use in this project).}
    \label{tab:repr-overview}
    \small
    \begin{tabular}{@{}p{2.6cm}p{4.2cm}p{5.8cm}@{}}
        \toprule
        \textbf{Family} & \textbf{Output shape / type} & \textbf{Role in the pipeline} \\
        \midrule
        BoW / TF-IDF & Sparse matrix ($N \times |\mathcal{V}|$) & Classical ML (NB, SVM, LR) with fast linear algebra \\
        FastText / AraVec & Dense word vectors averaged per doc & Bridging sparse lexicons and semantics \\
        AraBERT & Contextual vectors per token / \texttt{[CLS]} & Strong baseline for Arabic NLP \\
        \bottomrule
    \end{tabular}
\end{table}

\section{Bag of Words (BoW)}

\subsection{Principle}
The \textbf{bag-of-words} model ignores word order: each document is represented only by how often terms from $\mathcal{V}$ appear. After fixing an ordering of vocabulary terms $t_1,\ldots,t_{|\mathcal{V}|}$, document $d_i$ maps to
\[
    \mathbf{x}_i = [c_{i,1},\, c_{i,2},\, \ldots,\, c_{i,|\mathcal{V}|}]^{\top},
\]
where $c_{i,j}$ is the raw count (or binary presence) of term $t_j$ in $d_i$. In matrix form, the whole corpus is a document--term matrix $\mathbf{X} \in \mathbb{R}^{N \times |\mathcal{V}|}$, usually very sparse.

\subsection{\textit{n}-grams and vocabulary control}
Unigrams capture single tokens; bigrams and higher-order \textit{n}-grams retain short local context (\textit{e.g.}\ ``not good'' vs.\ ``good'') at the cost of a larger $|\mathcal{V}|$. In this codebase, \texttt{CountVectorizer} supports \texttt{ngram\_range=(1,2)} (unigrams + bigrams), \texttt{max\_features} to cap $|\mathcal{V}|$, and \texttt{min\_df} / \texttt{max\_df} to drop overly rare or overly frequent terms. For Arabic, \texttt{lowercase=False} is appropriate because case folding is irrelevant and can harm Latin fragments.

\subsection{Strengths and limits}
BoW is fast and interpretable, and pairs well with linear models. It cannot capture long-range dependencies or polysemy: the same surface form has a single dimension regardless of context. Morphological richness in Arabic inflates vocabulary size unless stemming/normalisation is applied.

\section{Term Frequency--Inverse Document Frequency (TF-IDF)}

\subsection{Motivation}
Raw counts emphasise long documents and frequent function words. \textbf{TF-IDF} reweights terms so that words discriminative for a few documents receive higher weight than words that appear everywhere.

\subsection{Components}
Let $f_{t,d}$ be the raw frequency of term $t$ in document $d$, and $\mathrm{df}(t)$ the number of documents in which $t$ appears. A standard weight is
\begin{equation}
    \mathrm{TF\text{-}IDF}(t,d) = w^{\mathrm{TF}}(t,d)\;\times\;w^{\mathrm{IDF}}(t),
    \label{eq:tfidf-product}
\end{equation}
with inverse document frequency often written as
\begin{equation}
    w^{\mathrm{IDF}}(t) = \log\frac{1 + N}{1 + \mathrm{df}(t)} + 1
    \label{eq:idf-smooth}
\end{equation}
(smoothed IDF as in scikit-learn), where $N$ is the corpus size. Term frequency $w^{\mathrm{TF}}$ may be raw counts, normalised counts, or---as in this project---\textbf{sublinear TF}: $w^{\mathrm{TF}} = 1 + \log f_{t,d}$ when $f_{t,d}>0$, which dampens the influence of very frequent repetitions within one document (\texttt{sublinear\_tf=True} in \texttt{TfidfVectorizer}).

\subsection{Normalisation}
The implementation applies \textbf{L2 normalisation} per row (\texttt{norm="l2"}), so each document vector lies on a hypersphere and cosine similarity equals the dot product---a common choice for text similarity and linear SVMs.

\subsection{Relation to BoW}
TF-IDF uses the same sparse indexing as BoW but replaces counts by TF-IDF weights. Feature selection (Chi-square, mutual information, etc.) can be applied on top of either representation.

\section{Distributed representations: FastText and AraVec}

\subsection{Word embeddings}
\textbf{FastText} (and comparable tools such as \textbf{AraVec}) learn dense vectors $\mathbf{e}_t \in \mathbb{R}^{k}$ for words (and subword units) from large unlabeled corpora using predictive objectives (skip-gram / CBOW). Semantically related words tend to have nearby vectors.

\subsection{Subwords and Arabic}
FastText represents each word as a sum of its character \textit{n}-grams. That helps \textbf{rare words}, \textbf{typos}, and morphological variants common in dialectal or noisy Arabic: the model can partially share evidence across forms that share character sequences.

\subsection{Document vectors}
For classification, a simple recipe is to average word vectors of tokens in the document (possibly weighted by TF-IDF). This yields a dense $\mathbf{x}_i \in \mathbb{R}^{k}$ suitable for logistic regression or SVM with linear kernel. Quality depends on whether the embedding vocabulary covers the domain (hotel reviews vs.\ tweets).

\section{Contextual representations: AraBERT}

\subsection{From words to transformers}
Unlike static embeddings, \textbf{AraBERT} (BERT architecture trained on Arabic text) builds \textbf{contextual} vectors: the representation of a token depends on the full sentence. A classification head often uses the pooled representation of the special \texttt{[CLS]} token or a mean pool over subword pieces produced by the WordPiece tokenizer.

\subsection{Fine-tuning}
Starting from pretrained weights, \textbf{fine-tuning} updates all (or top) layers on the downstream sentiment dataset with a task-specific softmax layer. This typically yields the best accuracy among the representations discussed here, at higher computational cost (GPU recommended).

\subsection{Arabic-specific benefit}
General multilingual models underperform compared to Arabic-centric pretraining when dialect, morphology, and spelling variation matter; AraBERT is designed for Modern Standard Arabic and much user-generated Arabic text used in benchmarks.

\section{Choosing a representation for this project}

\begin{itemize}[noitemsep]
    \item \textbf{Baseline experiments:} TF-IDF + linear SVM / logistic regression is a strong, reproducible baseline (see \texttt{tfidf()} defaults: \texttt{ngram\_range}, \texttt{max\_features}, \texttt{min\_df}, \texttt{max\_df}).
    \item \textbf{Feature selection studies:} Chi-square / MI / Fisher scores apply naturally on sparse BoW or TF-IDF matrices with varying \texttt{max\_features}.
    \item \textbf{Semantic gap:} FastText/AraVec averages help when exact word overlap between train and test is low.
    \item \textbf{Best accuracy budget:} Fine-tuned AraBERT is the usual ceiling model for Arabic sentiment, traded off against training time and inference cost.
\end{itemize}

% --- 6 ---
\chapter{Feature Selection}

High-dimensional sparse vectors from bag-of-words or TF---IDF can contain many weakly relevant dimensions. \textbf{Feature selection} keeps only the most informative dimensions (or keeps all as a baseline), which often \emph{reduces training time and memory} and can \emph{improve generalisation} when many noisy dimensions would otherwise dominate a linear decision boundary.

This project implements several classical selectors in \texttt{src/Feature\_Selection.py} and compares them in a reproducible way via \texttt{src/ml\_experiments.py} and the script \texttt{scripts/compare\_feature\_selection.py}.

\section{Notation}

Let $\mathbf{X}_{\mathrm{train}} \in \mathbb{R}^{n \times p}$ (typically sparse) be the matrix of vectorised training documents after a fixed feature extractor (e.g.\ TF---IDF), with binary labels $\mathbf{y} \in \{0,1\}^n$. Feature selection chooses a subset of $k \ll p$ column indices so that both train and test matrices are restricted to the same columns, avoiding leakage by fitting scores \emph{only} on the training fold.

\section{Methods implemented}

\subsection{Baseline: no selection (\texttt{none})}

\textbf{Idea:} use the full vectorised vocabulary (subject to the extractor's own \texttt{max\_features}, \texttt{min\_df}, etc.). This is the reference against which any selector must earn its keep.

\textbf{Implementation:} \texttt{select\_features(..., method="none")} returns \texttt{X\_train} and \texttt{X\_test} unchanged (\texttt{Feature\_Selection.py}).

\subsection{Chi-square test (\texttt{chi2})}

\textbf{Idea:} score each dimension (term or \textit{n}-gram) for \textbf{statistical dependence} with the class label under a $\chi^2$ test of independence, treating non-negative feature values as categorical counts in the spirit of discrete tests. Terms whose counts co-vary strongly with positive vs.\ negative documents receive high scores.

\textbf{Use case:} standard, fast filter for sparse \textbf{non-negative} text features (raw counts or TF---IDF with non-negative entries).

\textbf{Implementation:} \texttt{SelectKBest} with \texttt{score\_func=chi2}, keeping the top $k$ terms (\texttt{Select\_chi2}).

\subsection{Mutual information (\texttt{mi})}

\textbf{Idea:} \textbf{Mutual information} (MI) measures how much knowing a feature reduces uncertainty about the label (in bits, up to a constant). It can detect non-linear associations that a linear correlation might miss.

\textbf{Caveats:} MI can be slower and more memory-hungry than $\chi^2$ on very wide sparse matrices; scikit-learn may materialise dense blocks internally depending on configuration. It is still exposed as an optional filter for comparison.

\textbf{Implementation:} \texttt{SelectKBest} with \texttt{score\_func=mutual\_info\_classif} (\texttt{Select\_mi}).

\subsection{Fisher score ranking (\texttt{fisher})}

\textbf{Idea:} for each dimension $j$, compare \textbf{between-class spread} to \textbf{within-class variance}. A common scalar score is
\begin{equation}
    F_j = \frac{\sum_{c} n_c\, (\mu_{j,c} - \mu_j)^2}{\sum_{c} n_c\, \sigma_{j,c}^2 + \varepsilon},
\end{equation}
where $\mu_{j,c}$ and $\sigma_{j,c}^2$ are the mean and variance of feature $j$ in class $c$, $\mu_j$ is the global mean of feature $j$, $n_c$ is the class count, and $\varepsilon$ avoids division by zero. Larger $F_j$ suggests better class separation along dimension $j$.

\textbf{Implementation:} \texttt{Fisher\_score} computes per-dimension scores on the training matrix (dense if input was sparse); \texttt{Select\_fisher} keeps the indices of the top $k$ dimensions and slices train and test consistently.

\section{Selecting the best method: cross-validation over a discrete grid}

We treat the \textbf{choice of selector} as a \textbf{categorical hyperparameter} with a small grid
\[
    \mathcal{M} = \{\texttt{none},\, \texttt{chi2},\, \texttt{mi},\, \texttt{fisher}\}.
\]
This is analogous in spirit to \textbf{grid search}: instead of optimising real-valued weights, we \textbf{search over finitely many preprocessing choices} and pick the one with the best validation score.

\textbf{Important distinction:} scikit-learn's \texttt{GridSearchCV} is used elsewhere in the project for \textbf{numeric} hyperparameters (e.g.\ SVM $C$, vectoriser \texttt{max\_features}); the feature-selection comparison uses the same \textbf{evaluation protocol} (stratified $K$-fold on the training pool) but implements the loop explicitly so each fold \textbf{refits} the vectoriser, the selector, and the classifier only on training data---no test leakage.

\textbf{Procedure (summary):}
\begin{enumerate}[noitemsep]
    \item Split the corpus into a \textbf{training pool} and a \textbf{held-out test} set (stratified).
    \item For each method $m \in \mathcal{M}$: run \textbf{stratified $K$-fold CV} on the training pool. In each fold, vectorise train/validation, apply selection $m$ with fixed $k$, train a fixed classifier (e.g.\ linear SVM), record validation F$_1$.
    \item Choose $\arg\max_{m}$ of the \textbf{mean CV F$_1$} (tie-break with lower variance if desired).
    \item \textbf{Refit} the winning $m$ on the full training pool and report \textbf{hold-out} accuracy, precision, recall, and F$_1$ on the test set.
\end{enumerate}

\noindent
The helper \texttt{cross\_val\_score\_pipeline} in \texttt{ml\_experiments.py} encodes steps 2--3; \texttt{holdout\_evaluate} supports step~4. The CLI script \texttt{scripts/compare\_feature\_selection.py} prints a sorted table and the best method automatically.

\section{Code: running the comparison}

Listing~\ref{lst:fs-cli} shows the command-line entry point (run from the \textbf{repository root}). Listing~\ref{lst:fs-api} shows the equivalent programmatic pattern for a minimal notebook cell.

\begin{lstlisting}[style=py, caption={Shell: compare selectors with fixed TF---IDF and linear SVM (example).}, label={lst:fs-cli}]
python scripts/compare_feature_selection.py ^
  --dataset reviews ^
  --max-samples 12000 ^
  --extraction tfidf ^
  --methods none,chi2,mi,fisher ^
  --k-select 5000 ^
  --cv 5 ^
  --classifier linear_svc ^
  --output-csv report/feature_selection_cv.csv
\end{lstlisting}
\textit{Note:} On Unix shells, replace \texttt{\^{}} with line continuation \texttt{\textbackslash}.

\begin{lstlisting}[style=py, caption={Python: same logic inside a notebook or driver script.}, label={lst:fs-api}]
import sys
from pathlib import Path

ROOT = Path("..").resolve()   # adjust if your cwd differs
sys.path.insert(0, str(ROOT / "src"))

from ml_experiments import (
    init_rng,
    load_binary_xy,
    cross_val_score_pipeline,
    holdout_evaluate,
)
from sklearn.model_selection import train_test_split

SEED = 42
init_rng(SEED)
X, y = load_binary_xy("reviews", max_samples=12000, random_state=SEED)
X_train, X_test, y_train, y_test = train_test_split(
    X, y, test_size=0.2, random_state=SEED, stratify=y
)

mean_f1, std_f1 = cross_val_score_pipeline(
    X_train, y_train,
    extraction_method="tfidf",
    selection_method="chi2",
    k_select=5000,
    classifier_name="linear_svc",
    n_splits=5,
    random_state=SEED,
)
print("chi2:", mean_f1, "+/-", std_f1)
\end{lstlisting}

\section{Results (placeholder figure)}

Export a screenshot of the terminal or notebook table produced by the comparison (sorted mean F$_1$, standard deviation, and optional error column for failed runs). Save the image as \texttt{gs1.png} in the same directory as this \texttt{.tex} file (\texttt{report/gs1.png}) before compiling.

\begin{figure}[htbp]
    \centering
    \IfFileExists{gs1.png}{%
        \includegraphics[width=0.92\textwidth]{gs1.png}%
    }{%
        \fbox{\parbox[c][4.6cm][c]{0.92\textwidth}{\centering\textit{(Placeholder --- no file \texttt{gs1.png} yet)}\\[0.6em]%
        Export a screenshot of \texttt{compare\_feature\_selection.py} output and save it as \texttt{report/gs1.png}, then recompile.}}%
    }
    \caption{Feature selection comparison: cross-validated F$_1$ across methods (\texttt{none}, \texttt{chi2}, \texttt{mi}, \texttt{fisher}) with fixed extraction and classifier. Replace the placeholder by adding \texttt{gs1.png} next to this \texttt{.tex} file.}
    \label{fig:feature-selection-gs1}
\end{figure}

\section{Choosing $k$ (number of retained features)}

The hyperparameter $k$ in \texttt{SelectKBest} / Fisher truncation should be chosen on \textbf{validation} data or inner CV. If $k$ exceeds the number of columns produced by the vectoriser on a training fold, scikit-learn returns all available features and may emit a warning---reduce $k$ or increase \texttt{max\_features} in the extractor. Reporting a \textbf{curve} of F$_1$ vs.\ $k$ (e.g.\ $k \in \{500, 1000, 2000, 5000\}$) is a useful extension beyond the single-$k$ comparison table.

% --- 7 ---
\chapter{Models}

This chapter summarises the models used or considered in the project. The goal is not only to obtain a high score, but also to compare how classical sparse-feature classifiers and sequence-based deep models behave on Arabic tweets and hotel reviews.

\section{Machine Learning Models}

Classical machine learning models are trained on fixed-size vectors such as BoW or TF---IDF. In Arabic sentiment analysis, they are strong baselines because sentiment often depends on discriminative words and short phrases.

\subsection{Multinomial Naive Bayes}

\textbf{Definition.} Naive Bayes is a probabilistic classifier based on Bayes' rule and the assumption that features are conditionally independent given the class. For a document represented by features $\mathbf{x} = (x_1,\ldots,x_p)$, the predicted class is
\begin{equation}
    \hat{y}
    =
    \arg\max_{c}
    P(c)\prod_{j=1}^{p} P(x_j \mid c).
\end{equation}
In practice, logarithms are used:
\begin{equation}
    \hat{y}
    =
    \arg\max_{c}
    \left[
    \log P(c) + \sum_{j=1}^{p} x_j \log P(t_j \mid c)
    \right].
\end{equation}

\textbf{Suitability for this project.} Multinomial Naive Bayes is usually good as a fast baseline for BoW/TF---IDF text classification. It can perform well on short Arabic tweets when class-specific words are clear, but it may underperform linear SVM or logistic regression because the independence assumption is too simple for negation, dialectal expressions, and longer review sentences.

\subsection{Support Vector Machine (Linear SVM)}

\textbf{Definition.} A Support Vector Machine finds a separating hyperplane with the largest margin between classes. With a linear classifier, the prediction is
\begin{equation}
    \hat{y} = \mathrm{sign}(\mathbf{w}^{\top}\mathbf{x} + b).
\end{equation}
The soft-margin objective is commonly written as
\begin{equation}
    \min_{\mathbf{w},b,\xi}
    \frac{1}{2}\|\mathbf{w}\|^2
    + C\sum_{i=1}^{n}\xi_i,
    \quad
    y_i(\mathbf{w}^{\top}\mathbf{x}_i+b) \geq 1-\xi_i,
    \quad
    \xi_i \geq 0.
\end{equation}
The parameter $C$ controls the trade-off between a wide margin and training errors.

\textbf{Suitability for this project.} Linear SVM is usually one of the strongest classical models for Arabic sentiment with TF---IDF, especially in high-dimensional sparse spaces. It is well suited to the project because it handles many features efficiently and often gives robust performance on both tweets and reviews.

\subsection{Logistic Regression}

\textbf{Definition.} Logistic regression is a linear classifier that estimates the probability of the positive class. For binary sentiment classification,
\begin{equation}
    P(y=1 \mid \mathbf{x})
    =
    \sigma(\mathbf{w}^{\top}\mathbf{x}+b)
    =
    \frac{1}{1+\exp[-(\mathbf{w}^{\top}\mathbf{x}+b)]}.
\end{equation}
The predicted class is positive when this probability exceeds a threshold, usually $0.5$.

\textbf{Suitability for this project.} Logistic regression is a strong and interpretable baseline. It is usually competitive with linear SVM on TF---IDF features, and its learned weights can be inspected to identify which Arabic tokens or \textit{n}-grams push predictions toward positive or negative sentiment.

\section{Deep Learning Models}

Deep learning models can consume token sequences and learn internal representations instead of relying only on manually fixed sparse features. They can model local word patterns, word order, and longer dependencies, but require more training time and benefit from GPU acceleration.

\subsection{Convolutional Neural Network (CNN)}

\textbf{Definition.} A CNN applies filters over consecutive token embeddings to detect local patterns similar to learned \textit{n}-grams. If $\mathbf{e}_i$ is the embedding of token $i$, a convolutional filter of width $h$ produces features such as
\begin{equation}
    z_i = f(\mathbf{W}[\mathbf{e}_i;\mathbf{e}_{i+1};\ldots;\mathbf{e}_{i+h-1}] + b),
\end{equation}
where $f$ is a non-linear activation. A max-pooling layer then keeps the strongest detected pattern:
\begin{equation}
    z_{\max} = \max_i z_i.
\end{equation}

\textbf{Suitability for this project.} CNNs are often good for sentiment analysis because sentiment clues appear in short phrases. They are suitable for Arabic tweets and reviews when enough labeled data is available, but on smaller tweet subsets a TF---IDF + SVM baseline can still be stronger and more stable.

\subsection{Long Short-Term Memory (LSTM / BiLSTM)}

\textbf{Definition.} LSTM is a recurrent neural network designed to remember useful sequence information through gates. A simplified LSTM update is
\begin{align}
    \mathbf{i}_t &= \sigma(\mathbf{W}_i[\mathbf{h}_{t-1}, \mathbf{x}_t] + \mathbf{b}_i),\\
    \mathbf{f}_t &= \sigma(\mathbf{W}_f[\mathbf{h}_{t-1}, \mathbf{x}_t] + \mathbf{b}_f),\\
    \mathbf{o}_t &= \sigma(\mathbf{W}_o[\mathbf{h}_{t-1}, \mathbf{x}_t] + \mathbf{b}_o),\\
    \mathbf{c}_t &= \mathbf{f}_t \odot \mathbf{c}_{t-1} + \mathbf{i}_t \odot \tilde{\mathbf{c}}_t,\\
    \mathbf{h}_t &= \mathbf{o}_t \odot \tanh(\mathbf{c}_t).
\end{align}
A bidirectional LSTM (BiLSTM) reads the sequence from left to right and right to left, combining both contexts.

\textbf{Suitability for this project.} LSTM/BiLSTM models are useful for longer reviews where word order and negation can matter. However, they are more sensitive to dataset size, tokenisation, training epochs, and hardware. In this project they are useful as deep-learning baselines, but not guaranteed to outperform linear SVM on TF---IDF.

\subsection{AraBERT}

\textbf{Definition.} AraBERT is a transformer model pretrained on Arabic text. Transformers use self-attention, where each token representation attends to other tokens in the sequence. A simplified attention operation is
\begin{equation}
    \mathrm{Attention}(Q,K,V)
    =
    \mathrm{softmax}\left(\frac{QK^{\top}}{\sqrt{d_k}}\right)V.
\end{equation}
For classification, the model representation (often the \texttt{[CLS]} token or pooled output) is passed to a classification layer:
\begin{equation}
    \hat{\mathbf{y}} = \mathrm{softmax}(\mathbf{W}\mathbf{h}_{\texttt{[CLS]}}+\mathbf{b}).
\end{equation}

\textbf{Suitability for this project.} AraBERT is usually very good for Arabic sentiment analysis because it starts from language knowledge learned from large Arabic corpora. It is expected to outperform simple sparse-feature models when fine-tuned correctly, but it requires more memory, time, and preferably a GPU. For this project, AraBERT is best treated as an advanced model or future improvement if compute resources are limited.

% --- 8 ---
\chapter{Experimental Setup}

This chapter describes the practical setup used to run the experiments in the project. The same general protocol is followed for both datasets so that model comparisons are fair.

\section{Datasets and label preparation}

Two datasets are used:
\begin{itemize}[noitemsep]
    \item \textbf{Tweets dataset:} \texttt{ASTD/data/Tweets.txt}, originally labelled with \texttt{POS}, \texttt{NEG}, \texttt{NEUTRAL}, and \texttt{OBJ}. For binary experiments, only \texttt{POS} and \texttt{NEG} are kept and balanced.
    \item \textbf{Hotel reviews dataset:} \texttt{data/balanced-reviews.txt}, loaded with UTF-16 encoding. Star ratings are mapped as 4--5 $\rightarrow$ \texttt{POS} and 1--2 $\rightarrow$ \texttt{NEG}; rating 3 is dropped if present.
\end{itemize}

\section{Preprocessing}

Preprocessing is implemented in \texttt{src/preprocessing.py}. Tweets use a tweet-specific pipeline that handles URLs, mentions, hashtags, Arabic normalisation, elongation, punctuation, and optional emoji sentiment tokens. Reviews use a review-specific pipeline that additionally removes decorative quotes copied from table exports. Empty cleaned texts are removed before training.

\section{Data splitting and validation}

The standard split is:
\begin{itemize}[noitemsep]
    \item \textbf{80\% training pool} and \textbf{20\% held-out test set};
    \item \textbf{stratification} is used so positive and negative classes remain balanced in each split;
    \item method selection and hyperparameter search are performed using cross-validation on the training pool only;
    \item the held-out test set is used once for final reporting.
\end{itemize}
This avoids information leakage from the test set into feature selection or hyperparameter tuning.

\section{Feature extraction and feature selection}

Sparse text features are produced with:
\begin{itemize}[noitemsep]
    \item \textbf{BoW} using \texttt{CountVectorizer};
    \item \textbf{TF---IDF} using \texttt{TfidfVectorizer}, usually with unigrams and bigrams, \texttt{lowercase=False}, \texttt{min\_df}, \texttt{max\_df}, and \texttt{max\_features}.
\end{itemize}
Feature selection methods (\texttt{none}, \texttt{chi2}, \texttt{mi}, \texttt{fisher}) are compared with fixed extraction and classifier settings using \texttt{scripts/compare\_feature\_selection.py}. Feature extraction methods are compared with \texttt{scripts/compare\_feature\_extraction.py}.

\section{Models and hyperparameters}

The classical models are \textbf{Multinomial Naive Bayes}, \textbf{Linear SVM}, and \textbf{Logistic Regression}. Hyperparameters such as SVM $C$, TF---IDF vocabulary size, \textit{n}-gram range, and document-frequency thresholds are tuned with \texttt{scripts/gridsearch\_hyperparams.py}. Deep-learning experiments use Keras/TensorFlow CNN and BiLSTM notebooks with token sequences, embedding layers, a fixed maximum sequence length, and a small number of epochs for comparison.

\section{Software environment}

The project uses Python with \texttt{pandas}, \texttt{numpy}, \texttt{scikit-learn}, \texttt{matplotlib}, and \texttt{TensorFlow/Keras} for deep learning. Experiments are organised in reusable modules under \texttt{src/} and command-line scripts under \texttt{scripts/}, while notebooks are kept mostly as readable experiment reports.

% --- 9 ---
\chapter{Evaluation Metrics}

Evaluation metrics convert predictions into interpretable numbers. In binary sentiment classification, each prediction belongs to one of four cases:
\begin{itemize}[noitemsep]
    \item \textbf{TP} (true positive): a positive text predicted as positive;
    \item \textbf{TN} (true negative): a negative text predicted as negative;
    \item \textbf{FP} (false positive): a negative text predicted as positive;
    \item \textbf{FN} (false negative): a positive text predicted as negative.
\end{itemize}

\section{Confusion Matrix}

A \textbf{confusion matrix} is a table that counts correct and incorrect predictions for each class. For binary sentiment, using \texttt{NEG} as class 0 and \texttt{POS} as class 1, the matrix is
\[
\begin{array}{c|cc}
 & \textbf{Predicted NEG} & \textbf{Predicted POS} \\
\hline
\textbf{Actual NEG} & TN & FP \\
\textbf{Actual POS} & FN & TP
\end{array}
\]
This matrix is the starting point for accuracy, precision, recall, and F$_1$. It is also useful visually: a good model has large values on the diagonal ($TN$, $TP$) and small values off the diagonal ($FP$, $FN$).

\section{Accuracy}

Accuracy measures the proportion of all predictions that are correct:
\[
    \mathrm{Accuracy} = \frac{\mathrm{TP} + \mathrm{TN}}{\mathrm{TP} + \mathrm{TN} + \mathrm{FP} + \mathrm{FN}}.
\]
It is easy to understand, but it can be misleading if classes are imbalanced. For this reason, F$_1$ is also reported.

\section{Precision}

Precision answers: among the texts predicted as positive, how many were truly positive?
\[
    \mathrm{Precision} = \frac{\mathrm{TP}}{\mathrm{TP} + \mathrm{FP}}.
\]
High precision means few negative texts are incorrectly labelled as positive.

\section{Recall}

Recall answers: among all true positive texts, how many did the model find?
\[
    \mathrm{Recall} = \frac{\mathrm{TP}}{\mathrm{TP} + \mathrm{FN}}.
\]
High recall means the model misses few positive examples.

\section{F1-score}

The F$_1$-score is the harmonic mean of precision and recall:
\[
    F_1 = 2 \times \frac{\mathrm{Precision} \times \mathrm{Recall}}{\mathrm{Precision} + \mathrm{Recall}}.
\]
It is especially useful in this project because it balances false positives and false negatives better than accuracy alone.

\section{ROC / AUC}

ROC/AUC can be used when a model outputs scores or probabilities. The ROC curve plots true positive rate against false positive rate at different thresholds, and AUC summarises the curve in one value. In this project, AUC is optional because some strong models such as \texttt{LinearSVC} naturally output margins rather than calibrated probabilities.

% --- 10 ---
\chapter{Results and Discussion}
\textit{This is the core chapter.} Include tables comparing models and representations.

\begin{table}[htbp]
    \centering
    \caption{Example results (replace with your experiments)}
    \begin{tabular}{@{}llcc@{}}
        \toprule
        \textbf{Model} & \textbf{Representation} & \textbf{Accuracy} & \textbf{F1} \\
        \midrule
        SVM      & TF-IDF      & 91\% & 0.90 \\
        NB       & BoW         & 84\% & 0.82 \\
        AraBERT  & Transformer & 95\% & 0.94 \\
        \bottomrule
    \end{tabular}
\end{table}

Discuss:
\begin{itemize}[noitemsep]
    \item Which model performs best and why.
    \item HARD vs.\ ASTD differences.
    \item Impact of feature selection.
    \item Why stronger representations (e.g.\ AraBERT) may outperform sparse features.
\end{itemize}

Include figures: accuracy comparison, training time, effect of feature size.

% --- 11 ---
\chapter{Challenges}
Important for Arabic NLP:
\begin{itemize}[noitemsep]
    \item Dialects and code-switching.
    \item Slang and informal spelling.
    \item Ambiguity and negation.
    \item Noisy tweets vs.\ formal reviews.
    \item Preprocessing trade-offs.
\end{itemize}

% --- 12 ---
\chapter{Conclusion}
Summarize best methods, main findings, and lessons learned.

% --- 13 ---
\chapter{Possible Improvements and Future Work}
Examples:
\begin{itemize}[noitemsep]
    \item Larger or more diverse datasets.
    \item Multilingual or dialect-aware models.
    \item Dialect-specific preprocessing.
    \item Ensemble methods.
    \item Emotion detection beyond polarity.
\end{itemize}

% --- 14 ---
\begin{thebibliography}{99}
\bibitem{placeholder}
Add references: papers, dataset documentation, AraBERT, scikit-learn, etc.
\end{thebibliography}

\appendix
\chapter{Professional Report Checklist}
\begin{itemize}[noitemsep]
    \item Tables, confusion matrices, charts.
    \item Pipeline or architecture diagrams.
    \item Clear comparison and honest discussion of limitations.
\end{itemize}

\vspace{1em}
\noindent\textit{Structure derived from \texttt{docs/repport-struct.md}. Fill all bracketed and placeholder sections before submission.}

\end{document}
